\documentclass[a4paper,11pt]{ltjsarticle}
\usepackage[haranoaji]{luatexja-preset}
\usepackage{amsmath,amssymb}
\usepackage[top=12mm,bottom=10mm,left=14mm,right=14mm]{geometry}
\usepackage{array}
\usepackage{tikz}
\usetikzlibrary{calc}
\pagestyle{empty}
\setlength{\parindent}{0pt}
\newlength{\qcol}\setlength{\qcol}{0.47\textwidth}
\newlength{\qgap}
\newlength{\anscolw}\setlength{\anscolw}{0.30\textwidth}
\newlength{\ansh}
\newcommand{\Q}[2]{\parbox[t]{\qcol}{\raggedright (#1)\hspace{0.6em}$\displaystyle #2$}}
\newcommand{\QR}[4]{\Q{#1}{#2}\hfill\Q{#3}{#4}\par\vspace{\qgap}}
\newcommand{\anscell}[1]{\rule[-\ansdrop]{0pt}{\ansh}(#1)}
\newlength{\ansdrop}

\begin{document}
{\large\bfseries 計算問題　11}\hfill 名前(\underline{\hspace{32mm}})\par\vspace{3mm}
■（1）〜（16）の計算をしなさい。（17）〜（20）は方程式を解きなさい。\par\vspace{3mm}
\setlength{\qgap}{5.4mm}
\QR{1}{-10+(-13)}{2}{(-54)\div(+6)}
\QR{3}{(2a+7)+(3a-2)}{4}{\{(8-18)\div(-5)\}+6\times(-1)}
\QR{5}{-6\times4-48\div(-2^2)}{6}{\dfrac{-x+3}{2}\times8}
\QR{7}{-100a\div5}{8}{(+7)+(-5)+(-3)}
\QR{9}{(-14)\times(-7)}{10}{2(6a+4)-3(4a-12)}
\QR{11}{(24a-9)\div3}{12}{5-2\times4}
\QR{13}{4\div8\times(-16)}{14}{6+7^2+(-8)}
\QR{15}{\dfrac{1}{2}(4x+2)+\dfrac{1}{3}(6x-9)}{16}{5-(-6)\div\left(-\dfrac{3}{4}\right)}
\QR{17}{5x-8=-2x+6}{18}{\dfrac{1}{3}x+6=\dfrac{1}{2}x+4}
\QR{19}{0.03x-0.02=1}{20}{3(2x+5)=4(x-3)}
\vspace{1mm}
\Q{21}{1+2+3+\cdots+18+19+20}\par\vspace{3mm}
\setlength{\ansh}{9.5mm}\setlength{\ansdrop}{6mm}%
\begin{center}\begin{tabular}{|p{\anscolw}|p{\anscolw}|p{\anscolw}|}\hline
\anscell{1} & \anscell{2} & \anscell{3} \\\hline
\anscell{4} & \anscell{5} & \anscell{6} \\\hline
\anscell{7} & \anscell{8} & \anscell{9} \\\hline
\anscell{10} & \anscell{11} & \anscell{12} \\\hline
\anscell{13} & \anscell{14} & \anscell{15} \\\hline
\anscell{16} & \anscell{17} & \anscell{18} \\\hline
\anscell{19} & \anscell{20} & \anscell{21} \\\hline
\end{tabular}\end{center}

\end{document}
